\markdownRendererDocumentBegin
\markdownRendererWarning{The `hybrid` option has been soft-deprecated.}{The `hybrid` option has been soft-deprecated.}{Consider using one of the following better options for mixing TeX and markdown: `contentBlocks`, `rawAttribute`, `texComments`, `texMathDollars`, `texMathSingleBackslash`, and `texMathDoubleBackslash`. For more information, see the user manual at <https://witiko.github.io/markdown/>.}{Consider using one of the following better options for mixing TeX and markdown: `contentBlocks`, `rawAttribute`, `texComments`, `texMathDollars`, `texMathSingleBackslash`, and `texMathDoubleBackslash`. For more information, see the user manual at <https://witiko.github.io/markdown/>.}\markdownRendererSectionBegin
\markdownRendererHeadingOne{Introducción a la Programación Competitiva}\markdownRendererInterblockSeparator
{}\markdownRendererSectionBegin
\markdownRendererHeadingTwo{Veredictos}\markdownRendererInterblockSeparator
{}\begin{mdtable3}\markdownRendererSoftLineBreak
{}\rowcolor{headergray}\markdownRendererSoftLineBreak
{}\textbf{Código} & \textbf{Sinónimo} & \textbf{Posible explicación}\tabularnewline\markdownRendererSoftLineBreak
{}CE & Compilation Error & Error de sintaxis, faltan librerías o tipos mal definidos.\tabularnewline\markdownRendererSoftLineBreak
{}RTE & Run Time Error & Acceso a memoria inválida, división entre cero, etc.\tabularnewline\markdownRendererSoftLineBreak
{}TLE & Time Limit Exceeded & Algoritmo ineficiente o bucle infinito.\tabularnewline\markdownRendererSoftLineBreak
{}MLE & Memory Limit Exceeded & Uso excesivo de arrays o estructuras grandes.\tabularnewline\markdownRendererSoftLineBreak
{}WA & Wrong Answer & Lógica incorrecta o caso no cubierto.\tabularnewline\markdownRendererSoftLineBreak
{}OLE & Output Limit Exceeded & Imprime más líneas de lo esperado.\tabularnewline\markdownRendererSoftLineBreak
{}PE & Presentation Error & Saltos de línea o espacios extra.\tabularnewline\markdownRendererSoftLineBreak
{}AC & Accepted & El código fue aceptado.\tabularnewline\end{mdtable3}\markdownRendererInterblockSeparator
{}
\markdownRendererSectionEnd \markdownRendererSectionBegin
\markdownRendererHeadingTwo{Compilación}\markdownRendererInterblockSeparator
{}\begin{mdtable3}\markdownRendererSoftLineBreak
{}\rowcolor{headergray}\markdownRendererSoftLineBreak
{}\textbf{Sistema Operativo} & \textbf{Acción} & \textbf{Comando}\tabularnewline\markdownRendererSoftLineBreak
{}Linux/Mac OS & Compilar & g++ -o <problem> <problem>.cpp\tabularnewline\markdownRendererSoftLineBreak
{}Windows & Compilar & g++ -o <problem> <problem>.cpp\tabularnewline\markdownRendererSoftLineBreak
{}Linux/Mac OS & Ejecutar & <problem>\tabularnewline\markdownRendererSoftLineBreak
{}Windows & Ejecutar & <problem> .exe\tabularnewline\end{mdtable3}\markdownRendererParagraphSeparator
{}Para activar todas las advertencias o establecer una versión de C++ agrega los flags:\markdownRendererParagraphSeparator
{}\markdownRendererStrongEmphasis{-Wall -Wextra -std=c++<version>}:\markdownRendererInterblockSeparator
{}\markdownRendererInputFencedCode{./_markdown_main/a4bc4cb6cb013afe8c2f1ac5f7373f0e.verbatim}{cpp}{cpp}\markdownRendererInterblockSeparator
{}
\markdownRendererSectionEnd \markdownRendererSectionBegin
\markdownRendererHeadingTwo{Lectura}\markdownRendererInterblockSeparator
{}\markdownRendererSectionBegin
\markdownRendererHeadingThree{Strings con espacios}\markdownRendererInterblockSeparator
{}\markdownRendererInputFencedCode{./_markdown_main/ba6df3695673dfd3f1efeef9988309ff.verbatim}{cpp}{cpp}\markdownRendererInterblockSeparator
{}\markdownRendererSectionBegin
\markdownRendererHeadingFour{Limpiar buffer}\markdownRendererInterblockSeparator
{}Si se alterna con lectura de enteros, en ocasiones se guardara el último salto de línea, por lo que será necesario limpiar el buffer.\markdownRendererInterblockSeparator
{}\markdownRendererInputFencedCode{./_markdown_main/17137cea55572a9d7c3fe17a8c29a1cb.verbatim}{cpp}{cpp}\markdownRendererInterblockSeparator
{}
\markdownRendererSectionEnd 
\markdownRendererSectionEnd \markdownRendererSectionBegin
\markdownRendererHeadingThree{End of File}\markdownRendererInterblockSeparator
{}Si la cantidad de entrada es desconocida, leer hasta el final del archivo\markdownRendererInterblockSeparator
{}\markdownRendererInputFencedCode{./_markdown_main/c11393b1aa3f30c728313688ed568e61.verbatim}{cpp}{cpp}\markdownRendererInterblockSeparator
{}
\markdownRendererSectionEnd \markdownRendererSectionBegin
\markdownRendererHeadingThree{Archivos}\markdownRendererInterblockSeparator
{}Redirgir la entrada y salida a archivos\markdownRendererInterblockSeparator
{}\markdownRendererInputFencedCode{./_markdown_main/fa94a09e1291be07e2aaf8f4fea3d179.verbatim}{cpp}{cpp}
\markdownRendererSectionEnd 
\markdownRendererSectionEnd 
\markdownRendererSectionEnd \markdownRendererDocumentEnd